\section{The Assumption: Body and Soul in Heavenly Glory}

\definitionaxis{The Immaculate Mother of God, the ever-Virgin Mary, having completed the course of her earthly life, was assumed body and soul into heavenly glory.}{Pius XII, apostolic constitution \emph{Munificentissimus Deus}, 1 November 1950, definition at paragraph 44, after the 1946 episcopal consultation \emph{Deiparae Virginis Mariae}.}{The definition does not say whether Mary died, identify a date or place, canonize an apocryphal narrative, or equate her passive Assumption with Christ's Ascension by his own power.}

The Assumption is the Marian dogma of Easter. It does not add a second resurrection alongside Christ's. The risen Lord is ``firstfruits'' and cause; Mary's bodily glorification is an anticipated participation in the destiny he promises his members (1 Cor 15:20--23). The privilege is singular in timing and relation to Christ, while its substance---the redemption of the body and communion with God---is the hope of the whole Church.

\subsection{The exact economy of the definition}

Pius XII's subject is densely dogmatic: the one assumed is the Immaculate Mother of God and ever-Virgin. The four dogmas are thus gathered without being collapsed. ``Having completed the course of her earthly life'' marks the transition but deliberately avoids defining its physiological manner. ``Was assumed'' is passive: God acts. ``Body and soul'' excludes a merely spiritual survival or honor paid only to relics. ``Into heavenly glory'' names a transformed eschatological condition, not relocation of an unchanged organism to a point above the clouds.

Christ's Ascension and Mary's Assumption therefore differ essentially. The incarnate Son, risen by divine power, ascends as Lord and opens heaven (John 10:17--18; Acts 1:9--11; CCC 659--667). Mary is a redeemed creature received into glory by the power of her Son. Traditional shorthand can say that Christ ascends and Mary is assumed, provided it does not imply that Christ's humanity is independent of the Father and Spirit: the Paschal mystery is the work of the triune God.

\subsection{Scripture: resurrection promise and Marian fittingness}

The canonical Scriptures contain no narrative of Mary's final earthly moment or bodily Assumption. John Paul II states this directly in his general audience of 2 July 1997. The dogma is nevertheless not unscriptural. It belongs to the biblical revelation of Christ's victory, bodily resurrection, the destiny of the saints, and Mary's relation to him, as received in Tradition.

First Corinthians 15 is the primary doctrinal horizon. Christ is raised as firstfruits; at his coming those who belong to him will be raised; the perishable puts on imperishability; death is swallowed in victory (1 Cor 15:20--28, 42--57). Mary's Assumption anticipates but does not abolish this corporate sequence. Romans 8:11 promises that the Spirit who raised Jesus will give life to mortal bodies; 8:23 describes believers awaiting adoption's bodily consummation. Philippians 3:20--21 says Christ will transform the lowly body to conform it to his glorious body. These texts state the common hope directly. Mary's early participation is known through the Church's Tradition and dogmatic judgment.

Luke's Marian texts supply fittingness rather than a hidden end-of-life report. The one greeted as graced, blessed among women, and mother of the Lord confesses that the Mighty One has acted and that every generation will call her blessed (Luke 1:28, 42--49). The sword and discipleship lead through the Pasch; Christ's victory fittingly reaches the mother united to him by grace and flesh. Genesis 3:15 and the New Eve tradition place Mary within the complete enmity between the serpent and God's saving work. Since death and bodily corruption belong to sin's dominion, her bodily glory manifests the new Adam's victory.

Psalm 132:8---``Arise, O Lord, into your resting place, you and the ark of your might''---and Ark imagery have been applied typologically to the Assumption. Revelation 12's woman clothed with the sun supplies an ecclesial and Marian icon of glory and conflict. Neither text literally narrates Mary's final day. \emph{Munificentissimus Deus} 26--27 itself reports that theologians use such passages in accommodated and typological ways. A dogmatic exposition should show their place without upgrading them into modern historical reportage.

\subsection{Silence, apocrypha, and the first surviving narratives}

The New Testament's silence leaves room for history but forbids invented apostolic certainty. Late-antique \emph{Transitus Mariae} traditions narrate Mary's dormition, funeral, translation, or bodily glorification in differing forms. They are important evidence that Christians were meditating on her end and that a narrative complex developed. Their pseudonymous attributions, variations, miraculous elaborations, and disputed dating prevent their use as straightforward eyewitness depositions.

St.~Epiphanius of Salamis, writing in the fourth century, candidly says that Scripture does not disclose Mary's end and that no one knows it with certainty (\emph{Panarion} 78.11, within his discussion against the Antidicomarianites). This negative witness is historically consequential. It refutes claims of a universally explicit early narrative while not deciding what the Church could later recognize through Tradition. A mature history lets Epiphanius speak instead of forcing him into a fictitious Assumption catena.

The growth of a feast of Mary's commemoration or Dormition in the East, and its reception at Rome, is firmer evidence of ecclesial worship. Exact origin dates remain dependent on liturgical manuscripts and local histories, but by the early medieval period the feast is clearly established. \emph{Munificentissimus Deus} 16--20 notes the sacramentary sent by Pope Adrian I to Charlemagne, the stational procession associated with Pope Sergius I, and Pope Leo IV's addition of vigil and octave. Liturgy does not by itself provide a modern critical narrative of the event; sustained worship does reveal what the Church came to confess and celebrate.

\subsection{Mature patristic and Byzantine witness}

By the seventh and eighth centuries homilists state bodily glorification with clarity. Pseudo-Modestus's encomium speaks of Mary's resurrection from the tomb and reception into incorruption. St.~Germanus of Constantinople argues that the vessel of God could not remain under death's corruption and treats her continuing maternal relation to the Church. St.~John Damascene's Dormition homilies give the most influential synthesis: it was fitting that the one who preserved virginity in childbirth should have her body preserved from corruption, that the one who bore the Creator should dwell in the divine tabernacles, and that the Mother should possess what belongs to her Son (\emph{Second Homily on the Dormition} 14; cited in \emph{Munificentissimus Deus} 21).

These are substantial witnesses to the doctrine, but they are later than the apostolic age and often preach within the narrative world of Dormition traditions. Their historical role must be located accurately. The validity of their theological fittingness does not make each narrative detail dogma.

Eastern ``Dormition'' and Western ``Assumption'' should not be set against each other simplistically. Dormition emphasizes falling asleep or death; Assumption emphasizes bodily reception into glory. Much Eastern liturgy confesses both death and glorification. The Latin definition chose language broad enough not to bind the question of death. Ecumenical clarity requires honoring the Eastern tradition without falsely saying the 1950 formula defined every Dormition narrative.

\subsection{Did Mary die?}

Theologians have disputed whether Mary died, while the common liturgical and patristic tradition favors her Dormition. Pius XII's phrase ``having completed the course of her earthly life'' settles neither view. John Paul II's audience of 25 June 1997 presents Mary's death as the common tradition and theologically fitting participation in Christ's death, while explicitly acknowledging that the definition did not declare it a truth of faith.

The distinction between an open dogmatic question and a weighty traditional opinion is essential. One may argue that conformity to Christ, the liturgy of Dormition, and patristic consensus make death the better theological account. One may not add ``after she died'' to the defined proposition as though Pius XII had bound it. Conversely, one may not turn the definition's silence into a teaching that Mary did not die.

\subsection{Medieval theology and Aquinas's limited witness}

Medieval Western authors commonly receive Mary's bodily Assumption, even while the documentary basis and mode remain discussed. St.~Thomas never devotes a \emph{Summa} question to proving it. \emph{Munificentissimus Deus} 31 explicitly acknowledges this: whenever Thomas touches the subject, he holds that Mary's body as well as her soul is in heaven, but he never treats the question directly.

Pius XII's note identifies the \emph{Exposition of the Angelic Salutation}, the \emph{Exposition of the Apostles' Creed} article 5, and passages in the commentary on Book IV of the \emph{Sentences} (d.12, q.1, a.3, sol.3; d.43, q.1, a.3, sols.1--2). Their precise textual status and relation to Thomas's authenticated corpus should be respected. The safe claim is the one Pius XII makes: Thomistic witness is affirmative but occasional. It is false to invent a lost Thomistic demonstration or to make general resurrection metaphysics a direct historical proof.

Thomas nevertheless supplies illuminating principles. The soul is the substantial form of the body; human beatitude is not finally complete in permanent disembodiment (ST I, q.76; Supplement, qq.75--86). Christ's Resurrection is efficient and exemplary cause of ours (ST III, q.56, a.1). Grace perfects nature, and glory consummates grace. Applied under the dogma's authority, these principles show why Mary's body matters: salvation fulfills the human person rather than discarding embodied identity.

\subsection{From feast and consensus to definition}

Papal teaching increasingly articulates the Assumption before 1950. Pius XII's constitution surveys feast, devotion, Fathers, theologians, and predecessors rather than claiming a new revelation. In \emph{Deiparae Virginis Mariae} (1946), he asked the bishops whether they judged the bodily Assumption revealable and definable and whether they desired definition. Their responses were almost unanimously affirmative (\emph{Munificentissimus Deus} 12).

Episcopal consultation was not a plebiscite that created truth. The bishops were witnesses to the faith of their churches and participants in the ordinary teaching office. Pius XII regarded their moral unanimity, together with the historical evidence, as a certain sign that the doctrine was contained in the deposit entrusted to the Church. The later solemn act identified that revealed truth definitively.

On 1 November 1950, invoking the authority of Christ, Peter and Paul, and his own, Pius XII pronounced, declared, and defined ``as a divinely revealed dogma'' that the Immaculate Mother of God, ever-Virgin Mary, having completed her earthly life, was assumed body and soul into heavenly glory (\emph{Munificentissimus Deus} 44). The formula's economy should be quoted exactly when legal or dogmatic force matters and paraphrased without additions elsewhere.

\subsection{Ecclesial and anthropological meaning}

The Assumption is a promise made visible. In Mary the Church contemplates the beginning and image of what she will be in the age to come (LG 68; CCC 966, 972). The doctrine therefore opposes both spiritualism and body-idolatry. The body is neither a disposable shell nor self-sufficient; the whole person is created, redeemed, and destined for transformed communion. Mary's female body is not honored because it escaped creaturehood but because grace brings creaturehood to glory.

Her glory remains ecclesial rather than private possession. She is not removed into indifference toward pilgrims. Within the communion of saints she praises God and intercedes; every such act depends upon the Lamb, the one heavenly High Priest. Nor does bodily glory make her omnipresent, omniscient by nature, or a divine distributor of grace. Created participation remains created even at its summit.

\subsection{Errors excluded}

The Assumption is not Christ's Ascension repeated by a second saving subject. It is not a doctrine that Mary's body was never material, that she did not age, or that Christians may neglect burial and bodily works of mercy. It does not certify a claimed tomb, house, date, or geographic tradition. It does not define her death or deny it. It does not turn Revelation 12 into literal video, or apocryphal \emph{Transitus} texts into apostolic memoir. Above all, it does not remove Mary from the universal redemption: the Immaculate one is assumed because Christ has conquered, and her anticipated glory magnifies the firstfruits rather than competing with him.
